% BHU PYQ -- Manually transcribed & error-corrected
\documentclass[11pt,a4paper]{article}

\usepackage[a4paper,top=2.5cm,bottom=2.5cm,left=2.8cm,right=2.8cm,headheight=14pt]{geometry}
\usepackage{amsmath,amssymb}
\usepackage[T1]{fontenc}
\usepackage[utf8]{inputenc}
\usepackage{lmodern}
\usepackage{microtype}
\usepackage{fancyhdr}
\usepackage{enumitem}
\usepackage{hyperref}
\usepackage{lastpage}
\usepackage{parskip}

\hypersetup{colorlinks=true,urlcolor=black,linkcolor=black,
  pdftitle={BPT-401: Electromagnetic Theory and Basic Electronics (Old Course) -- BHU PYQ},pdfauthor={Banaras Hindu University}}

\pagestyle{fancy}\fancyhf{}
\renewcommand{\headrulewidth}{0.4pt}\renewcommand{\footrulewidth}{0.4pt}
\fancyhead[L]{\small   \textit{Banaras Hindu University}}
\fancyhead[R]{\small   \textit{BPT-401: EMT and Basic Electronics (Old)}}
\fancyfoot[C]{\small Page  \thepage\ of \pageref{LastPage}}


\newlist{parts}{enumerate}{2}
\setlist[parts,1]{label=(\alph*),leftmargin=*,itemsep=5pt,topsep=3pt}
\setlist[parts,2]{label=(\roman*),leftmargin=*,itemsep=3pt,topsep=2pt}

\newcommand{\pts}[1]{\hfill   \textit{[#1]}}


\begin{document}


\begin{center}
  {\large\bfseries BANARAS HINDU UNIVERSITY}\\[2pt]
  {\normalsize B.Sc. (Hons.) Semester IV Examination, 2013-14 (Old Course)}\\[6pt]
  \rule{0.8\linewidth}{0.5pt}\\[6pt]
  {\large\bfseries Paper No. BPT-401: Electromagnetic Theory and Basic Electronics}\\[4pt]
  {\normalsize Time: 3 Hours\hfill Maximum Marks: 70}
\end{center}
\rule{\linewidth}{0.4pt}
\smallskip



\noindent   \textit{Instructions: Answer   \textbf{five} questions including Question~1 which is compulsory. Symbols have their usual meanings.}
\bigskip




\noindent   \textbf{Question 1.} Answer any   \textbf{five} of the following: \pts{5 $ \times$ 4 = 20}

\begin{parts}
  \item Prove that the divergence of the curl of any vector field vanishes: $
abla \cdot (
abla  \times \vec{A}) = 0$.
  \item Define scalar and vector potentials. How do you express the electric field $\vec{E}$ and magnetic field $\vec{B}$ in terms of these potentials?
  \item Write down the four Maxwell's equations in differential form and briefly explain the physical meaning of each.
  \item A material has conductivity $\sigma = 10^{-3}\, \text{S/m}$ and relative permittivity $\varepsilon_r = 2$. At what frequency would the conduction current density be equal to the displacement current density?
  \item Describe the position of the Fermi level in an intrinsic and an extrinsic semiconductor. Discuss the effect of temperature on its position.
  \item Why is a capacitor filter mostly used in a power supply rather than an inductor filter? Explain.
  \item Explain how the current gain $\alpha$ is different from $\beta$. Establish the relationship between $\alpha$ and $\beta$ for a transistor.
\end{parts}

\medskip\rule{\linewidth}{0.2pt}




\noindent   \textbf{Question 2.}

\begin{parts}
  \item State and prove Gauss's divergence theorem. \pts{5}
  \item Derive the electromagnetic wave equation for the propagation of a plane wave in a non-conducting isotropic dielectric medium. Obtain expressions for the propagation constant, intrinsic impedance, phase velocity, and group velocity. \pts{9}
\end{parts}

\medskip\rule{\linewidth}{0.2pt}




\noindent   \textbf{Question 3.}

\begin{parts}
  \item State and prove Poynting's theorem for the conservation of energy in an electromagnetic field, and discuss the physical significance of each term in the equation. \pts{8}
  \item Define depth of penetration (skin depth). Show that electromagnetic waves are heavily attenuated inside a good conductor. \pts{6}
\end{parts}
\hfill   \textbf{OR}

\begin{parts}
  \item A plane wave travelling in a homogenous, isotropic, and linear dielectric medium is incident obliquely at a plane interface of another dielectric. Determine the reflection and transmission coefficients if the incident wave is polarized with its electric vector lying perpendicular to the plane of incidence. \pts{9}
  \item What is a transmission line? Explain the different parameters of a transmission line. Determine the characteristic impedance of a lossless cable if the inductance per unit length $L = 0.25\,\mu \text{H/m}$ and capacitance per unit length $C = 100\, \text{pF/m}$. \pts{5}
\end{parts}

\medskip\rule{\linewidth}{0.2pt}




\noindent   \textbf{Question 4.}

\begin{parts}
  \item Explain the depletion region in a p-n junction. Show that the width of the depletion region in the p-n junction decreases with the increase of doping impurity concentration. \pts{8}
  \item Describe the working of a single-stage RC coupled common-emitter amplifier. Plot the frequency response of the amplifier and explain its features. \pts{6}
\end{parts}
\hfill   \textbf{OR}

\begin{parts}
  \item Explain the working of a full-wave bridge rectifier having a capacitor filter with a suitable circuit diagram. Obtain the relationship for the ripple factor. \pts{8}
  \item Why does a transistor need biasing? Draw the circuit diagram of a transistor self-bias (voltage divider bias) circuit. Explain how it provides better thermal stability than a fixed-bias circuit. \pts{6}
\end{parts}

\medskip\rule{\linewidth}{0.2pt}




\noindent   \textbf{Question 5.}

\begin{parts}
  \item Starting from Maxwell's equations, derive the wave equation for the propagation of a plane electromagnetic wave in free space. Obtain the solution representing a plane wave propagating along the $+z$ direction. \pts{7}
  \item Find the Brewster angle for a parallel-polarized wave travelling from air into glass for which $\varepsilon_r = 5.0$. \pts{4}
  \item Explain the meaning of the Early effect. How does it affect the collector current in a junction transistor? \pts{3}
\end{parts}

\vfill
\rule{\linewidth}{0.4pt}

\begin{center}\small
  \href{https://bhu-exam.vercel.app/}{Click here to visit our website}\\[4pt]
  \href{https://me-aryan.vercel.app/}{Click here to visit our contributor}\\[4pt]
  \href{https://quashan-me.vercel.app/}{Click here to visit our contributor}
\end{center}

\end{document}
